\documentclass[11pt,a4paper]{article}
\usepackage[T1]{fontenc}
\usepackage[margin=28mm]{geometry}
\usepackage{amsmath,amssymb,amsthm,booktabs}
\usepackage[hidelinks]{hyperref}
\hypersetup{pdfauthor={Qichao Liu},pdftitle={The least denominator offset in the Erdos-Straus equation on prime squares}}
\urlstyle{same}
\newtheorem{theorem}{Theorem}
\newtheorem{lemma}[theorem]{Lemma}
\newtheorem{proposition}[theorem]{Proposition}
\newtheorem{corollary}[theorem]{Corollary}
\theoremstyle{definition}
\newtheorem{example}{Example}
\newcommand{\jac}[2]{\left(\frac{#1}{#2}\right)}
\newcommand{\ES}{Erd\H{o}s--Straus}

\title{The least denominator offset in the \ES{} equation on prime squares}
\author{Qichao Liu}
\date{26 November 2026}

\begin{document}
\maketitle

\begin{abstract}
For an odd prime \(p\), let \(T(p)\) be the least value of \(4x-p^2\)
among the ordered representations \(4/p^2=1/x+1/y+1/z\) with
\(x\le y\le z\). We determine the limiting distribution of \(T(p)\) over
the primes. Its atoms, in increasing order, are the primes congruent to
\(3\bmod 4\) and the integers \(3q\) with \(q\) prime and congruent to
\(1\bmod 4\). The \(j\)-th atom has mass \(2^{-j}\). For each fixed
offset, a Jacobi symbol gives a necessary condition that is sufficient
for a set of relative density one in each admissible residue class.
For every prime offset \(r\equiv3\bmod4\), we also construct an
arithmetic progression on which an explicit representation attains
\(T(p)=r\).
\end{abstract}

\noindent\textit{Keywords.} Unit fractions, \ES{} equation, prime squares,
quadratic reciprocity, relative natural density.

\section{The question and the main theorem}

The \ES{} conjecture asserts that for every integer \(n\ge2\) there are
positive integers \(x\le y\le z\) satisfying
\[
  \frac4n=\frac1x+\frac1y+\frac1z.
\]
Since \(1/x<4/n\), the smallest denominator must exceed \(n/4\).
For \(n=p^2\) with \(p\) an odd prime, the possible positive offsets
\(t=4x-p^2\) are \(3,7,11,15,\ldots\).
Call \(t\) \emph{feasible} for \(p\) if \(x=(p^2+t)/4\) is the smallest
denominator in such a representation. Write \(T(p)\) for the first
feasible offset and set \(T(p)=\infty\) if there is no representation.
For a set \(S\) of odd primes, its relative natural density is
\[
  \delta(S)=\lim_{X\to\infty}
  \frac{\#\{p\in S:p\le X\}}{\#\{p\le X:p\text{ is an odd prime}\}},
\]
when the limit exists.

\begin{theorem}\label{thm:law}
List the members of
\[
  \{q:q\ \text{prime},\ q\equiv3\pmod4\}
  \ \cup\
  \{3q:q\ \text{prime},\ q\equiv1\pmod4\}
\]
in increasing order as \(u_1,u_2,\ldots\).
For every \(j\ge1\),
\(\delta(\{p:T(p)=u_j\})=2^{-j}\).
Every other finite value has density zero, and so does \(T(p)=\infty\).
\end{theorem}

For example, \(3,7,11,15\) have masses \(1/2,1/4,1/8,1/16\),
respectively. The offset \(27\) has mass zero because its quadratic
character already occurs at \(3\). The table gives the counts for the
\(1{,}217\) primes in \(41\le p<10{,}000\). For each prime we applied
Lemma~\ref{lem:div} to successive offsets through \(31\). The last
column counts primes for which every tested offset failed.

\begin{center}
\small
\begin{tabular}{lrrrrrrrrr}
\toprule
First offset & 3 & 7 & 11 & 15 & 19 & 23 & 27 & 31 & \(>31\)\\
\midrule
Number of primes & 611 & 309 & 154 & 76 & 31 & 24 & 0 & 9 & 3\\
\bottomrule
\end{tabular}
\end{center}

\section{A local criterion at a fixed offset}

Fix \(t>0\) with \(t\equiv3\bmod4\). Throughout this section
\(p\nmid t\) and \(p^2>t\), conditions which exclude only finitely
many primes for fixed \(t\). Put
\[
  x=\frac{p^2+t}{4},\qquad M=p^2x.
\]
Then \(\gcd(M,t)=1\), since \(4x\equiv p^2\bmod t\).
After the term \(1/x\) has been removed, the remaining fraction is
\(t/M\).

\begin{lemma}\label{lem:div}
The offset \(t\) is feasible if and only if \(M^2\) has a positive
divisor \(d\equiv-M\pmod t\).
\end{lemma}

\begin{proof}
The identity \(t/M=1/y+1/z\) gives
\[
  (ty-M)(tz-M)=M^2.
\]
Both factors are positive since \(y,z>M/t\). Conversely, if \(d\)
is as stated, put \(e=M^2/d\). As \(M\) is invertible modulo \(t\),
we have \(e\equiv-M\bmod t\). Thus
\[
  y=(M+d)/t,\qquad z=(M+e)/t
\]
are positive integers satisfying the identity. Both exceed
\(M/t=xp^2/t>x\), so we may order them to obtain the asserted
representation.
\end{proof}

Write \(\chi_t(v)=\jac{v}{t}\) for the Jacobi symbol.

\begin{proposition}\label{prop:fixed}
If \(t\) is feasible, then \(\chi_t(p)=-1\). Conversely, for each
invertible residue class \(a\bmod t\) satisfying \(\chi_t(a)=-1\),
the proportion of primes in that class for which \(t\) is feasible
tends to one.
\end{proposition}

\begin{proof}
Let \(q\) be an odd prime dividing \(x\). Then
\(p^2\equiv-t\bmod q\), and \(q\) is coprime to \(pt\).
Quadratic reciprocity and \(t\equiv3\bmod4\) give
\[
  \jac{q}{t}=\jac{-t}{q}=1.
\]
If \(2\mid x\), then \(t=4x-p^2\equiv7\bmod8\), so
\(\chi_t(2)=1\) as well. Every prime factor of \(x\) has character
\(+1\). If \(\chi_t(p)=1\), the same holds for every divisor of
\(M^2\), whereas \(\chi_t(-M)=-1\) since \(\chi_t(-1)=-1\).
Lemma~\ref{lem:div} rules out feasibility.

For the converse, fix \(a\) with \(\chi_t(a)=-1\) and define
\(b\bmod t\) by \(4b\equiv-a^3\bmod t\). Then \(b\) is invertible
and \(\chi_t(b)=1\). For every sufficiently large prime
\(q\equiv b\bmod t\), quadratic reciprocity yields
\(\jac{-t}{q}=\jac{q}{t}=1\). The congruence \(p^2\equiv-t\bmod q\)
therefore has two distinct nonzero roots modulo \(q\).
If \(p\equiv a\bmod t\) belongs to either class, then \(q\mid x\).
In that case \(d=pq\) divides \(M^2\), and
\[
  4d\equiv4ab\equiv-a^4\equiv-4M\pmod t .
\]
Thus \(t\) is feasible by Lemma~\ref{lem:div}.

For a fixed finite set \(S\) of these auxiliary primes \(q\), the
Chinese remainder theorem and the prime number theorem for arithmetic
progressions give
\[
  \prod_{q\in S}\left(1-\frac{2}{q-1}\right)
\]
as the limiting proportion, among primes \(p\equiv a\bmod t\), that
avoid both roots for every \(q\in S\). This product bounds the upper
relative density of the primes where \(t\) is infeasible. The sum of
the reciprocals of the primes \(q\equiv b\bmod t\) diverges, so the
finite products can be made arbitrarily small.
\end{proof}

\section{Distribution of the first feasible offset}

\begin{proof}[Proof of Theorem~\ref{thm:law}]
Fix an integer \(B\ge3\). Proposition~\ref{prop:fixed} shows that on
a set of primes of density one, feasibility agrees with
\(\chi_t(p)=-1\) simultaneously for all \(t\le B\) with
\(t>0\) and \(t\equiv3\bmod4\). Indeed, only finitely many offsets
occur, and for each one there are only finitely many residue classes.
We may also discard the finitely many primes dividing any such \(t\)
or satisfying \(p^2\le B\).

Let \(Q(B)\) be the set of prime factors occurring in these offsets.
Then
\[
  Q(B)=\{q\le B:q\equiv3\pmod4\}
       \ \cup\ \{q\le B/3:q\equiv1\pmod4\},
\]
where \(q\) ranges over primes. A prime \(q\equiv3\bmod4\) first
occurs at \(t=q\). If \(q\equiv1\bmod4\) divides an offset \(t\),
then \(t/q\equiv3\bmod4\), so its first occurrence is at \(t=3q\).

Now every \(\chi_t(p)\) up to \(B\) equals \(+1\) if and only if
\(\jac{p}{q}=1\) for each \(q\in Q(B)\). One direction follows
from multiplicativity of the Jacobi symbol in its denominator.
For the other, take \(t=q\) whenever \(q\equiv3\bmod4\). For
\(q\equiv1\bmod4\) in \(Q(B)\), take \(t=3q\) and use
\(\jac{p}{3}=1\), already known from \(t=3\).

Exactly half of the nonzero classes modulo each \(q\) are squares.
The Chinese remainder theorem makes the choices for different \(q\)
independent, and the prime number theorem for arithmetic progressions
gives equal limiting frequency to all the resulting invertible
classes. With \(m(B)=\#Q(B)\) we obtain
\[
  \delta(\{p:T(p)>B\})=2^{-m(B)}.
\]
This formula also holds for \(B=2\), with \(Q(2)\) empty.
When \(B\) reaches \(u_j\), one new prime enters \(Q(B)\), so
the tail density drops by \(2^{-(j-1)}-2^{-j}=2^{-j}\).
At every other integer \(B\) it has no drop. Finally the upper
density of \(T(p)=\infty\) is at most \(2^{-j}\) for every \(j\).
\end{proof}

\paragraph{Relation to earlier work.}
Yamamoto used quadratic characters to exclude squares from his
congruence covering \cite[Theorem~2]{yamamoto}.
Elsholtz and Tao proved that odd squares admit no solutions of their
Types~I and~II \cite[Proposition~1.6]{elsholtztao}.
Lemma~\ref{lem:div} tests a prescribed smallest denominator across
all solution types. Geometric weights already appear in
Erd\H{o}s's distribution of the least quadratic nonresidue
\cite[equation~(7)]{erdos}. Here the density-one converse in
Proposition~\ref{prop:fixed} passes from the distribution of
quadratic characters to the distribution of actual representations.

\section{Arithmetic progressions with a prescribed first offset}

For the prime atoms of Theorem~\ref{thm:law}, an explicit construction
attains the first offset.

\begin{theorem}\label{thm:construction}
For every prime \(r\equiv3\bmod4\) there exists an invertible residue
class such that every prime \(p>r\) in that class satisfies \(T(p)=r\).
Consequently, for any integer \(H\ge1\) there are infinitely many
soluble inputs \(p^2\) for which none of the first \(H\) integers
above \(p^2/4\) can be the smallest denominator.
\end{theorem}

\begin{proof}
Let \(L\) be the least common multiple of all positive
\(t<r\) with \(t\equiv3\bmod4\), taking \(L=1\) if there are none.
Since \(r\) is prime and coprime to \(2L\), the conditions
\[
  p\equiv1\pmod{2L},\qquad p\equiv-1\pmod r
\]
specify one invertible residue class modulo \(2Lr\).
Dirichlet's theorem gives infinitely many primes in it. For \(p>r\)
and every earlier offset \(t\), we have \(\chi_t(p)=1\), so
Proposition~\ref{prop:fixed} excludes \(t\).
At \(t=r\), set \(x=(p^2+r)/4\) and \(k=(p+1)/r\). Then \(rk=p+1\)
gives the ordered representation
\[
  \frac4{p^2}=\frac1x+\frac1{pxk}+\frac1{p^2xk}.
\]
Hence \(T(p)=r\). Choosing a prime \(r\equiv3\bmod4\) with
\(r>4H-1\) excludes all of the first \(H\) candidates.
\end{proof}

\begin{corollary}\label{cor:quant}
There is an absolute constant \(c>0\) such that for infinitely many
soluble \(n=p^2\), the least feasible first denominator \(x\) satisfies
\[
  x-\frac n4\ge c\log n .
\]
\end{corollary}

\begin{proof}
Let \(r\) run through the primes congruent to \(3\bmod4\), and choose
the least prime \(p\) in the residue class modulo \(D=2Lr\) of
Theorem~\ref{thm:construction}. Since \(p\) is odd and
\(p\equiv-1\bmod r\), we have \(p\ge2r-1\). Linnik's theorem gives
absolute constants \(C,A>0\) such that \(p\le CD^A\)
\cite{heathbrown}. Now \(Lr\) divides \(\operatorname{lcm}(1,\ldots,r)\),
whose logarithm is \(O(r)\) by Chebyshev's estimate. Thus
\(\log n=2\log p=O(r)\), whereas Theorem~\ref{thm:construction}
gives \(x-n/4=r/4\). The primes \(p\) obtained in this way are
unbounded.
\end{proof}

\section{Examples and scope}

\begin{example}
Take \(r=11\), so \(L=21\). The prime \(p=43\) satisfies both
\(p\equiv1\bmod42\) and \(p\equiv-1\bmod11\). For \(n=1849\)
the first candidates \(463,464,465\) correspond to \(t=3,7,11\).
The first two offsets fail because \(43\equiv1\) modulo both \(3\)
and \(7\). The third gives \(k=4\) and
\[
  \frac4{1849}=\frac1{465}+\frac1{79980}+\frac1{3439140}.
\]
\end{example}

\begin{example}
Take \(p=29\) and \(t=19\), so \(x=215=5\cdot43\). Although
\(\jac{29}{19}=-1\), this offset fails. Modulo \(19\), the element
\(g=10\) has order \(18\), since \(g^9=-1\) and \(g^6=11\).
We have \(29\equiv g\) and \(5\equiv43\equiv g^2\), so
\(M=29^2\cdot5\cdot43\equiv g^6\). The target \(-M\) is \(g^{15}\).
Every divisor of \(M^2\) is congruent to \(g^{i+2j}\) for some
\(0\le i,j\le4\). These exponents are at most \(12\), and none is
\(15\bmod18\). Lemma~\ref{lem:div} excludes the offset. This gives
an exception to pointwise sufficiency in Proposition~\ref{prop:fixed}.
\end{example}

The results sample inputs \(n=p^2\). A proof of the \ES{} conjecture
for all prime inputs would imply it for all integers \(n\ge2\) by
scaling the denominators \cite{elsholtztao}. Our analysis uses that
\(-t\) is a square modulo odd primes dividing \((p^2+t)/4\).
For prime inputs, Dahan studies search depth and gives an obstruction
at the fixed offset \(7\) \cite{dahan}. The general conjecture remains
open.

\paragraph{Manuscript status and use of AI.}
OpenAI Codex proposed the problem, developed the statements and proofs,
searched for related literature, drafted the manuscript and wrote the
code for the finite numerical checks.
This version has not undergone
independent peer review. The literature comparison above identifies
related results. Priority for the distribution theorem has not been
established.

\begin{thebibliography}{9}

\bibitem{yamamoto}
K.~Yamamoto,
On the Diophantine equation \(4/n=1/x+1/y+1/z\),
\textit{Memoirs of the Faculty of Science, Kyushu University,
Series~A} \textbf{19} (1965), 37--47.
\href{https://doi.org/10.2206/kyushumfs.19.37}{doi:10.2206/kyushumfs.19.37}.

\bibitem{elsholtztao}
C.~Elsholtz and T.~Tao,
Counting the number of solutions to the \ES{} equation on unit fractions,
\textit{Journal of the Australian Mathematical Society}
\textbf{94} (2013), 50--105.
\href{https://doi.org/10.1017/S1446788712000468}{doi:10.1017/S1446788712000468}.

\bibitem{erdos}
P.~Erd\H{o}s,
Sz\'amelm\'eleti megjegyz\'esek, I,
\textit{Matematikai Lapok} \textbf{12} (1961), 10--17.
\href{https://www.renyi.hu/~p_erdos/1961-23.pdf}{Original paper}.

\bibitem{heathbrown}
D.~R. Heath-Brown,
Zero-free regions for Dirichlet \(L\)-functions, and the least prime
in an arithmetic progression,
\textit{Proceedings of the London Mathematical Society} (3)
\textbf{64} (1992), 265--338.
\href{https://doi.org/10.1112/plms/s3-64.2.265}{doi:10.1112/plms/s3-64.2.265}.

\bibitem{dahan}
B.~Dahan,
Sieve dimension and search depth for the \ES{} conjecture,
\(n\equiv1\pmod{24}\), preprint, version~2 (2026).
\href{https://arxiv.org/abs/2608.24035v2}{arXiv:2608.24035v2}.

\end{thebibliography}
\end{document}
